\chapter{Обзор методов анализа многоканальных нейрофизиологических данных}
\label{ch:literature_review}

%\textit{Основные результаты, изложенные в данной главе, опубликованы в работах соискателя %\cite{vakbib1, confbib1, scbib1}.}

\section{Модель генерации ЭЭГ/МЭГ сигналов и нейрональные ритмы}
\label{sec:eeg_model}

Одна из целей анализа многоканальных нейрофизиологических данных заключается в выявлении и интерпретации компонент, отражающих текущее функциональное состояние мозга в зависимости от экспериментального условия. В рамках линейной квазистатической модели распространения ЭЭГ/МЭГ сигнала, наблюдаемая сенсорная активность $\mathbf{X} \in \mathbb{R}^{N_c \times N_t}$ (где $N_c$ --- количество каналов, $N_t$ --- количество отсчетов времени) представляется как линейная суперпозиция вкладов распределённых источников:
\begin{equation}
	\mathbf{X} = \mathbf{A}\mathbf{S} + \boldsymbol{\eta},
\end{equation}
где $\mathbf{A} \in \mathbb{R}^{N_c \times N_s}$ --- матрица прямой модели (оператор смешивания), $\mathbf{S} \in \mathbb{R}^{N_s \times N_t}$ --- временные ряды в пространстве латентных источников, а $\boldsymbol{\eta}$ --- шумовая компонента.

Для решения обратной задачи на практике ищут матрицу пространственных фильтров $\mathbf{W} \in \mathbb{R}^{N_c \times N_s}$, позволяющую получить оценку сигналов источников $\hat{\mathbf{S}} = \mathbf{W}^T \mathbf{X}$ \cite{parra2005recipes}.

\section{Линейные методы пространственной фильтрации}
\label{sec:linear_methods}

В рамках описанной модели предложен широкий класс методов:
\begin{itemize}
	\item \textbf{Метод главных компонент (PCA)} максимизирует дисперсию проекций.
	\item \textbf{Анализ независимых компонент (ICA)} минимизирует взаимную информацию между выделенными компонентами \cite{kim2006iva}.
	\item \textbf{Канонический корреляционный анализ (CCA)} решает задачу поиска пар фильтров для двух наборов данных $\mathbf{X}$ и $\mathbf{Y}$ с целью максимизации корреляции $\rho = \text{corr}(\mathbf{w}_x^T\mathbf{X}, \mathbf{w}_y^T\mathbf{Y})$. Существуют модификации для многопробных ЭЭГ-данных, такие как TRCA \cite{tanaka2013trca} и CORRCA \cite{zhang2018corrca, parra2019corrca}.
\end{itemize}

Огромным преимуществом линейных методов является их физическая интерпретируемость: они позволяют преобразовать веса фильтра $\mathbf{w}$ в паттерны активации $\mathbf{a}$ (колонки матрицы $\mathbf{A}$) согласно формуле \cite{haufe2014interpretation}:
\begin{equation}
	\mathbf{a} = \frac{\mathbf{\Sigma}_{X}\mathbf{w}}{\mathbf{w}^T\mathbf{\Sigma}_{X}\mathbf{w}},
\end{equation}
где $\mathbf{\Sigma}_{X}$ --- ковариационная матрица сенсорных данных.

\section{Риманова геометрия пространственных ковариационных матриц}
\label{sec:riemannian_geometry}

В задачах декодирования ЭЭГ (например, моторного воображения) матрица ковариации $\mathbf{\Sigma} = \frac{1}{N_t-1}\mathbf{X}\mathbf{X}^T$ содержит ключевую пространственную информацию о распределении мощности источников. Однако ковариационные матрицы принадлежат не евклидову пространству, а нелинейному многообразию симметричных положительно определенных (SPD) матриц $\mathcal{S}_{++}^{N_c}$ \cite{barachant2013classification}.

Использование стандартных евклидовых метрик для таких матриц приводит к искажениям (например, эффект <<набухания>> --- swelling effect). Для корректной работы алгоритмов машинного обучения матрицы проецируются в касательное пространство (Tangent Space) в базовой точке (как правило, геометрическом среднем Фреше $\mathbf{\Sigma}_{ref}$), используя логарифмическое отображение \cite{yger2017riemann, suh2021riemannian}:
\begin{equation}
	\mathbf{S}_i = \log \left( \mathbf{\Sigma}_{ref}^{-1/2} \mathbf{\Sigma}_i \mathbf{\Sigma}_{ref}^{-1/2} \right),
\end{equation}
где $\mathbf{S}_i$ --- симметричная матрица в касательном пространстве, элементы которой можно векторизовать (получив евклидов вектор признаков) без потери топологической информации \cite{xu2020tangent}.

\section{Выделение индуцированной активности (семейство SPoC)}
\label{sec:induced_spoc}

Индуцированная активность связана с событием по времени, но обладает высокой вариабельностью фазы между пробами. При наличии внешней поведенческой переменной $z$ задача решается методом SPoC \cite{dahne2014spoc}, максимизирующим ковариацию между мощностью сигнала и переменной $z$:
\begin{equation}
	\mathbf{w}_{opt} = \arg\max_{\mathbf{w}} \frac{\mathbf{w}^T \mathbf{C}_z \mathbf{w}}{\mathbf{w}^T \mathbf{C} \mathbf{w}}, \quad \text{где } \mathbf{C}_z = \frac{1}{N_e} \sum_{e=1}^{N_e} (z_e - \bar{z}) \mathbf{C}_e.
\end{equation}
Алгоритм mSPoC \cite{dahne2013mspoc} обобщает этот подход для многомерных поведенческих переменных $\mathbf{Z} \in \mathbb{R}^{d \times N_e}$, одновременно оптимизируя фильтры $\mathbf{w}$ для ЭЭГ и весовые векторы $\mathbf{v}$ для таргета, аналогично концепции CCA.

\section{Нелинейное снижение размерности (UMAP)}
\label{sec:nonlinear_methods}

Методы нелинейного вложения, такие как Isomap \cite{tenenbaum2000isomap} и Laplacian Eigenmaps \cite{belkin2003laplacian}, строят низкоразмерные представления, опираясь на графы ближайших соседей. Современным стандартом является алгоритм UMAP \cite{mcinnes2018umap}, который минимизирует нечеткую кросс-энтропию между высокомерным и низкоразмерным графовыми представлениями данных. 

Для систематизации методов в таблице~\ref{tab:methods_comparison} приведено их сравнение. Как видно из таблицы, нелинейные методы хорошо выявляют структуру, но лишены интерпретируемости (нет обратного оператора), тогда как линейные методы (SPoC) нуждаются в явном задании таргета. Данное противоречие формирует пробел, на решение которого направлен предлагаемый в диссертации метод.

\begin{table}[htbp]
	\caption{Сравнение методов выделения компонент в нейрофизиологических данных}
	\label{tab:methods_comparison}
	\centering
	\footnotesize 
	\begin{tabular}{|p{2.1cm}|p{1.9cm}|p{3.5cm}|p{2.5cm}|p{2.3cm}|p{2.2cm}|}
		\hline
		\textbf{Метод} & \textbf{Тип} & \textbf{Оптимизируемый критерий} & \textbf{Использование внеш. переменных} & \textbf{Целевая активность} & \textbf{Интерпрет. (паттерн)} \\ \hline
		\textbf{PCA} & Линейный & Максимизация дисперсии & Нет & Вызванная & Да \\ \hline
		\textbf{ICA} & Линейный & Максимизация негауссовскости & Нет & Вызванная, артефакты & Да \\ \hline
		\textbf{UMAP} & Нелинейный & Кросс-энтропия топологических графов & Нет / Supervised & Любая (динамика) & Нет \\ \hline
		\textbf{SPoC} & Линейный & Отношение Рэлея для 1D таргета & Да (1D переменная) & Индуцированная & Да \\ \hline
		\textbf{mSPoC} & Линейный & Комодуляция мощности и nD таргета & Да (nD переменная) & Индуцированная & Да \\ \hline
		\textbf{TriCo (наш)} & Гибридный & Комодуляция мощности с геометрией манифолда & Да (координаты UMAP как таргет) & Индуцированная & \textbf{Да} \\ \hline
	\end{tabular}
\end{table}